\documentclass[10pt,aspectratio=169]{beamer}
\usepackage{bma02}

\title[La integral definida]{La integral definida}
\subtitle{Partición, refinamiento, integral inferior y superior.\\Área debajo de una parábola\\[2pt]\small Semana 2 -- Sesión 1}

\begin{document}

\begin{frame}
  \titlepage
\end{frame}

\begin{frame}{Objetivos de la sesión}
  \begin{enumerate}
    \item Definir \resaltar{partición} de un intervalo, su norma y sus refinamientos.
    \item Construir las \resaltar{sumas inferior y superior} de Darboux y estudiar cómo cambian al refinar.
    \item Definir integral inferior, integral superior y \resaltar{función integrable}.
    \item Calcular, desde la definición, el área bajo la parábola $y=x^2$.
  \end{enumerate}
  \vspace{1mm}
  {\small\textbf{Competencia 1 del sílabo:} Evalúa interpretaciones y propiedades del cálculo integral y su utilidad en problemas de la vida real.}
  \begin{block}{Plan (2 horas; incluye verificación previa de 5 min)}
    \begin{tabular}{ll}
      0:00 -- 0:15 & El problema del área: idea de aproximar con rectángulos\\
      0:15 -- 0:55 & Particiones, sumas inferior/superior, refinamiento\\
      0:55 -- 1:15 & Integral inferior y superior. Integrabilidad. Un contraejemplo\\
      1:15 -- 1:45 & Área debajo de una parábola\\
      1:45 -- 2:00 & Actividad colaborativa, resumen y cierre metacognitivo
    \end{tabular}
  \end{block}
\end{frame}

\begin{frame}{Aula invertida: preparación y verificación}
  \small
  \begin{previa}[antes de la sesión]
    \begin{enumerate}
      \item Lea las definiciones de partición, sumas inferior y superior e integrabilidad.
      \item Para $f(x)=x^2$ en $[0,1]$ calcule $L$ y $U$ con $n=2$.
      \item Pregunta guía: ¿cuándo coinciden las integrales inferior y superior?
    \end{enumerate}
  \end{previa}
  \begin{block}{Verificación al inicio (5 min, en pares)}
    \begin{enumerate}
      \item ¿Qué es un refinamiento de una partición?
      \item ¿Cómo se relacionan $L(f,P)$ y $U(f,P)$?
    \end{enumerate}
    Los pares discuten y votan; el docente aclara solo lo que genera dudas y ajusta el ritmo de la sesión.
  \end{block}
\end{frame}

\section{El problema del área}

\begin{frame}{¿Qué es el área bajo una curva?}
  \begin{columns}
    \begin{column}{0.5\textwidth}
      Para polígonos sabemos calcular áreas. Para una región limitada por una curva,
      \resaltar{aproximamos} con rectángulos:
      \begin{itemize}
        \item rectángulos \textbf{inscritos} (por debajo) $\Rightarrow$ aproximación \emph{por defecto};
        \item rectángulos \textbf{circunscritos} (por encima) $\Rightarrow$ aproximación \emph{por exceso}.
      \end{itemize}
      \begin{idea}
        Si al hacer los rectángulos cada vez más angostos ambas aproximaciones se acercan al mismo
        número, ese número es \emph{el área}.
      \end{idea}
    \end{column}
    \begin{column}{0.5\textwidth}
      \centering
      \begin{tikzpicture}
        \begin{axis}[bma, xmin=0, xmax=2.3, ymin=0, ymax=4.6, width=6.8cm, height=5.6cm,
            xtick={0,2}, ytick={4}]
          \addplot[fill=uniGranate!20, draw=uniGranate, ybar interval, samples=5, domain=0:2] {(x+0.5)^2};
          \addplot[fill=uniAzul!40, draw=uniAzul, ybar interval, samples=5, domain=0:2] {x^2};
          \addplot[curva, domain=0:2.1, samples=100] {x^2} node[pos=1, left] {$y=x^2$};
        \end{axis}
      \end{tikzpicture}

      \small Azul: suma inferior. Granate: exceso de la suma superior.
    \end{column}
  \end{columns}
\end{frame}

\section{Particiones y sumas de Darboux}

\begin{frame}{Partición de un intervalo}
  \begin{definicion}[partición]
    Una \resaltar{partición} de $[a,b]$ es un conjunto finito
    \[ P=\{x_0,x_1,\dots,x_n\}\quad\text{con}\quad a=x_0<x_1<\dots<x_n=b. \]
    Los subintervalos son $I_i=[x_{i-1},x_i]$, de longitud $\Delta x_i=x_i-x_{i-1}$.
    La \resaltar{norma} de $P$ es $\|P\|=\max\limits_{1\le i\le n}\Delta x_i$.
  \end{definicion}
  \begin{center}
    \begin{tikzpicture}[scale=1.2]
      \draw[thick] (0,0) -- (9,0);
      \foreach \x/\l in {0/x_0=a,1.3/x_1,2.1/x_2,3.8/x_3,5.0/x_4,6.9/x_5,9/x_6=b} {
        \draw[thick] (\x,0.12) -- (\x,-0.12) node[below] {$\l$};
      }
      \draw[decorate,decoration={brace,amplitude=5pt},uniGranate] (5.0,0.2) -- (6.9,0.2)
        node[midway,above=5pt] {$\Delta x_5=\|P\|$};
    \end{tikzpicture}
  \end{center}
  \begin{ejemplo}
    \begin{itemize}
      \item $P=\{0,\tfrac14,\tfrac12,1\}$ es partición de $[0,1]$ con $\|P\|=\tfrac12$.
      \item \textbf{Partición regular} en $n$ partes: $x_i=a+i\dfrac{b-a}{n}$, \ $\Delta x_i=\dfrac{b-a}{n}=\|P\|$.
    \end{itemize}
  \end{ejemplo}
\end{frame}

\begin{frame}{Refinamiento}
  \begin{definicion}[refinamiento]
    Si $P$ y $P^*$ son particiones de $[a,b]$ y $P\subset P^*$, decimos que $P^*$ es un
    \resaltar{refinamiento} de $P$ (o que $P^*$ es \emph{más fina} que $P$).
  \end{definicion}
  \begin{center}
    \begin{tikzpicture}[scale=1.1]
      \draw[thick] (0,0) -- (8,0); \node[left] at (-0.2,0) {$P$};
      \foreach \x in {0,2,5,8} \draw[thick] (\x,0.12) -- (\x,-0.12);
      \draw[thick] (0,-1) -- (8,-1); \node[left] at (-0.2,-1) {$P^*$};
      \foreach \x in {0,2,5,8} \draw[thick] (\x,-0.88) -- (\x,-1.12);
      \foreach \x in {1,3.5,6.5} \draw[thick,uniGranate] (\x,-0.85) -- (\x,-1.15);
      \node[uniGranate,right] at (8.2,-1) {\small puntos agregados};
    \end{tikzpicture}
  \end{center}
  \begin{observacion}
    \begin{itemize}
      \item Dadas dos particiones $P_1$, $P_2$, la unión $P_1\cup P_2$ es un refinamiento común de ambas.
      \item Al refinar, la norma no aumenta: $\|P^*\|\le\|P\|$.
    \end{itemize}
  \end{observacion}
\end{frame}

\begin{frame}{Sumas inferior y superior}
  Sea $f$ \resaltar{acotada} en $[a,b]$ y $P$ una partición. En cada $I_i$ definimos
  \[ m_i=\inf_{x\in I_i}f(x),\qquad M_i=\sup_{x\in I_i}f(x). \]
  \begin{definicion}[sumas de Darboux]
    \[ L(f,P)=\sum_{i=1}^n m_i\,\Delta x_i\quad\text{(suma inferior)},\qquad
       U(f,P)=\sum_{i=1}^n M_i\,\Delta x_i\quad\text{(suma superior)}. \]
  \end{definicion}
  \begin{columns}
    \begin{column}{0.48\textwidth}
      \centering
      \begin{tikzpicture}
        \begin{axis}[bma, xmin=0, xmax=3.3, ymin=0, ymax=2.6, width=6.4cm, height=4.3cm, xtick=\empty, ytick=\empty]
          \addplot[fill=uniAzul!35, draw=uniAzul] coordinates {(0.3,0) (0.3,1.95) (0.9,1.95) (0.9,0)};
          \addplot[fill=uniAzul!35, draw=uniAzul] coordinates {(0.9,0) (0.9,1.61) (1.5,1.61) (1.5,0)};
          \addplot[fill=uniAzul!35, draw=uniAzul] coordinates {(1.5,0) (1.5,0.8) (2.1,0.8) (2.1,0)};
          \addplot[fill=uniAzul!35, draw=uniAzul] coordinates {(2.1,0) (2.1,0.7) (2.7,0.7) (2.7,0)};
          \addplot[curva, domain=0.2:2.9] {1.5+0.8*sin(deg(2*x))};
        \end{axis}
      \end{tikzpicture}\\
      \small $L(f,P)$: altura $=$ mínimo en cada $I_i$
    \end{column}
    \begin{column}{0.48\textwidth}
      \centering
      \begin{tikzpicture}
        \begin{axis}[bma, xmin=0, xmax=3.3, ymin=0, ymax=2.6, width=6.4cm, height=4.3cm, xtick=\empty, ytick=\empty]
          \addplot[fill=uniGranate!30, draw=uniGranate] coordinates {(0.3,0) (0.3,2.3) (0.9,2.3) (0.9,0)};
          \addplot[fill=uniGranate!30, draw=uniGranate] coordinates {(0.9,0) (0.9,2.28) (1.5,2.28) (1.5,0)};
          \addplot[fill=uniGranate!30, draw=uniGranate] coordinates {(1.5,0) (1.5,1.61) (2.1,1.61) (2.1,0)};
          \addplot[fill=uniGranate!30, draw=uniGranate] coordinates {(2.1,0) (2.1,0.88) (2.7,0.88) (2.7,0)};
          \addplot[curva, domain=0.2:2.9] {1.5+0.8*sin(deg(2*x))};
        \end{axis}
      \end{tikzpicture}\\
      \small $U(f,P)$: altura $=$ máximo en cada $I_i$
    \end{column}
  \end{columns}
\end{frame}

\begin{frame}{Primeras propiedades}
  Si $m\le f(x)\le M$ en $[a,b]$, como $m\le m_i\le M_i\le M$:
  \[ \boxed{m(b-a)\le L(f,P)\le U(f,P)\le M(b-a)} \]
  \begin{teorema}[Lema del refinamiento]
    Si $P^*$ es un refinamiento de $P$, entonces
    \[ L(f,P)\le L(f,P^*)\qquad\text{y}\qquad U(f,P^*)\le U(f,P). \]
    Refinar \emph{mejora} ambas aproximaciones.
  \end{teorema}
  \textbf{Idea de la prueba.} Basta agregar un punto $y\in(x_{k-1},x_k)$. Los ínfimos en
  $[x_{k-1},y]$ y $[y,x_k]$ son $\ge m_k$, así que
  \[ m_k\,\Delta x_k\le\Big(\inf_{[x_{k-1},y]}f\Big)(y-x_{k-1})+\Big(\inf_{[y,x_k]}f\Big)(x_k-y). \]
\end{frame}

\begin{frame}{Visualizando el refinamiento}
  \centering
  \begin{tikzpicture}
    \begin{axis}[bma, name=a1, xmin=0, xmax=2.2, ymin=0, ymax=4.5, width=5.4cm, height=4.4cm, xtick={0,2}, ytick=\empty, title={\small $n=2$}]
      \addplot[fill=uniGranate!20, draw=uniGranate, ybar interval, samples=3, domain=0:2] {(x+1)^2};
      \addplot[fill=uniAzul!40, draw=uniAzul, ybar interval, samples=3, domain=0:2] {x^2};
      \addplot[curva, domain=0:2] {x^2};
    \end{axis}
    \begin{axis}[bma, at={(a1.east)}, anchor=west, xshift=0.6cm, name=a2, xmin=0, xmax=2.2, ymin=0, ymax=4.5, width=5.4cm, height=4.4cm, xtick={0,2}, ytick=\empty, title={\small $n=4$}]
      \addplot[fill=uniGranate!20, draw=uniGranate, ybar interval, samples=5, domain=0:2] {(x+0.5)^2};
      \addplot[fill=uniAzul!40, draw=uniAzul, ybar interval, samples=5, domain=0:2] {x^2};
      \addplot[curva, domain=0:2] {x^2};
    \end{axis}
    \begin{axis}[bma, at={(a2.east)}, anchor=west, xshift=0.6cm, xmin=0, xmax=2.2, ymin=0, ymax=4.5, width=5.4cm, height=4.4cm, xtick={0,2}, ytick=\empty, title={\small $n=16$}]
      \addplot[fill=uniGranate!20, draw=uniGranate, ybar interval, samples=17, domain=0:2] {(x+0.125)^2};
      \addplot[fill=uniAzul!40, draw=uniAzul, ybar interval, samples=17, domain=0:2] {x^2};
      \addplot[curva, domain=0:2] {x^2};
    \end{axis}
  \end{tikzpicture}

  \vspace{2mm}
  \begin{tabular}{c|ccc}
    $n$ & $L(f,P_n)$ & $U(f,P_n)$ & $U-L$\\ \hline
    2 & 1.000 & 5.000 & 4.000\\
    4 & 1.750 & 3.750 & 2.000\\
    16 & 2.4219 & 2.9219 & 0.500\\
    100 & 2.6268 & 2.7068 & 0.080
  \end{tabular}
  \quad
  \begin{minipage}{6cm}
    \small $f(x)=x^2$ en $[0,2]$. La parte granate (diferencia $U-L$) se hace cada vez más pequeña.
    ¿Hacia qué número convergen $L$ y $U$?
  \end{minipage}
\end{frame}

\begin{frame}{Toda suma inferior es menor que toda suma superior}
  \begin{teorema}
    Para cualesquiera particiones $P_1,P_2$ de $[a,b]$:\qquad $L(f,P_1)\le U(f,P_2)$.
  \end{teorema}
  \textbf{Prueba.} Sea $P^*=P_1\cup P_2$ (refinamiento común). Por el lema:
  \[ L(f,P_1)\le L(f,P^*)\le U(f,P^*)\le U(f,P_2).\qquad\blacksquare \]
  \begin{center}
    \begin{tikzpicture}[scale=1.1]
      \draw[thick,->] (0,0) -- (10,0);
      \foreach \x in {0.6,1.4,2.3,3.0,3.6} \fill[uniAzul] (\x,0) circle (2pt);
      \foreach \x in {6.2,6.8,7.6,8.5,9.3} \fill[uniGranate] (\x,0) circle (2pt);
      \node[uniAzul,above] at (2,0.15) {sumas inferiores $L(f,P)$};
      \node[uniGranate,above] at (7.8,0.15) {sumas superiores $U(f,P)$};
      \draw[uniAzul,thick] (4.4,-0.25) -- (4.4,0.25) node[above] {$\underline{\int}$};
      \draw[uniGranate,thick] (5.4,-0.25) -- (5.4,0.25) node[above] {$\overline{\int}$};
      \node[below] at (4.9,-0.3) {\small ¿coinciden?};
    \end{tikzpicture}
  \end{center}
  El conjunto de sumas inferiores está acotado superiormente (por cualquier $U$) y el de superiores
  acotado inferiormente (por cualquier $L$).
\end{frame}

\section{Integral inferior, superior e integrabilidad}

\begin{frame}{Integral inferior y superior}
  \begin{definicion}
    Sea $f$ acotada en $[a,b]$.
    \[ \underline{\int_a^b}f=\sup_{P}\,L(f,P)\quad\text{(integral inferior)},\qquad
       \overline{\int_a^b}f=\inf_{P}\,U(f,P)\quad\text{(integral superior)}. \]
    Siempre se cumple $\displaystyle\underline{\int_a^b}f\le\overline{\int_a^b}f$.
  \end{definicion}
  \begin{definicion}[integral de Riemann]
    $f$ es \resaltar{integrable} en $[a,b]$ si $\displaystyle\underline{\int_a^b}f=\overline{\int_a^b}f$.
    El valor común se denota $\displaystyle\int_a^b f(x)\dx$.
  \end{definicion}
  \begin{observacion}
    Si $f\ge0$, $\int_a^bf(x)\dx$ \emph{define} el área de la región bajo la gráfica de $f$.
  \end{observacion}
\end{frame}

\begin{frame}{Criterio de integrabilidad de Riemann}
  \begin{teorema}[Criterio de Riemann]
    $f$ acotada es integrable en $[a,b]$ si y solo si para todo $\varepsilon>0$ existe una partición $P$ tal que
    \[ U(f,P)-L(f,P)<\varepsilon. \]
  \end{teorema}
  \begin{columns}
    \begin{column}{0.55\textwidth}
      $U(f,P)-L(f,P)=\displaystyle\sum_{i=1}^n(M_i-m_i)\Delta x_i$ es el área de los ``rectángulos de
      oscilación'' (en granate en la figura). Integrable $\iff$ esa área puede hacerse tan pequeña como se quiera.
    \end{column}
    \begin{column}{0.45\textwidth}
      \centering
      \begin{tikzpicture}
        \begin{axis}[bma, xmin=0, xmax=2.2, ymin=0, ymax=4.5, width=5.8cm, height=4.3cm, xtick={0,2}, ytick=\empty]
          \addplot[fill=uniGranate!40, draw=uniGranate, ybar interval, samples=9, domain=0:2] {(x+0.25)^2};
          \addplot[fill=white, draw=uniGranate, ybar interval, samples=9, domain=0:2] {x^2};
          \addplot[curva, domain=0:2] {x^2};
        \end{axis}
      \end{tikzpicture}
    \end{column}
  \end{columns}
\end{frame}

\begin{frame}{Una función que no es integrable}
  \begin{ejemplo}[función de Dirichlet]
    \[ f(x)=\begin{cases}1,&x\in\Q\\0,&x\notin\Q\end{cases}\qquad\text{en }[0,1]. \]
    Todo subintervalo contiene racionales e irracionales, luego $m_i=0$ y $M_i=1$ para toda partición:
    \[ L(f,P)=0,\qquad U(f,P)=\sum\Delta x_i=1. \]
    Por tanto $\underline{\int_0^1}f=0\ne1=\overline{\int_0^1}f$: \ $f$ \alert{no es integrable}.
  \end{ejemplo}
  \begin{center}
    \begin{tikzpicture}
      \begin{axis}[bma, xmin=0, xmax=1.1, ymin=-0.2, ymax=1.3, width=8cm, height=3.4cm, xtick={1}, ytick={1}]
        \addplot[only marks, mark=*, mark size=0.6pt, uniGranate, samples=120, domain=0:1] {1};
        \addplot[only marks, mark=*, mark size=0.6pt, uniAzul, samples=120, domain=0.004:1.004] {0};
      \end{axis}
    \end{tikzpicture}
  \end{center}
\end{frame}

\section{Área debajo de una parábola}

\begin{frame}{Área bajo $y=x^2$ en $[0,b]$: planteamiento}
  Partición regular: $x_i=\dfrac{ib}{n}$, \ $\Delta x=\dfrac bn$. Como $f(x)=x^2$ es creciente en $[0,b]$:
  \[ m_i=f(x_{i-1})=\frac{(i-1)^2b^2}{n^2},\qquad M_i=f(x_i)=\frac{i^2b^2}{n^2}. \]
  \begin{block}{Fórmula que necesitamos}
    \[ \sum_{i=1}^n i^2=\frac{n(n+1)(2n+1)}{6} \qquad\text{(se prueba por inducción).} \]
  \end{block}
  Entonces
  \[ U(f,P_n)=\sum_{i=1}^n\frac{i^2b^2}{n^2}\cdot\frac bn=\frac{b^3}{n^3}\cdot\frac{n(n+1)(2n+1)}{6}
    =\frac{b^3}{6}\left(1+\frac1n\right)\left(2+\frac1n\right). \]
\end{frame}

\begin{frame}{Área bajo $y=x^2$ en $[0,b]$: conclusión}
  Análogamente, con $\sum_{i=1}^{n}(i-1)^2=\frac{(n-1)n(2n-1)}{6}$:
  \[ L(f,P_n)=\frac{b^3}{6}\left(1-\frac1n\right)\left(2-\frac1n\right). \]
  Por lo tanto, para todo $n$:
  \[ L(f,P_n)\le\underline{\int_0^b}x^2\dx\le\overline{\int_0^b}x^2\dx\le U(f,P_n). \]
  Haciendo $n\to\infty$, ambos extremos tienden a $\dfrac{b^3}{3}$. Luego
  \begin{teorema}[Arquímedes, $\sim$250 a.C.]
    \[ \int_0^b x^2\dx=\frac{b^3}{3}. \]
    El área del segmento parabólico es $\tfrac13$ del rectángulo $b\times b^2$ que lo contiene.
  \end{teorema}
  \small Compare con la tabla: con $b=2$, $\frac83\approx2.6667$. \checkmark
\end{frame}

\begin{frame}{Interpretación: un tercio del rectángulo}
  \begin{columns}
    \begin{column}{0.5\textwidth}
      \centering
      \begin{tikzpicture}
        \begin{axis}[bma, xmin=0, xmax=2.4, ymin=0, ymax=4.6, width=6.5cm, height=5.6cm, xtick={2}, xticklabels={$b$}, ytick={4}, yticklabels={$b^2$}]
          \addplot[draw=none, fill=uniGranate!15] coordinates {(0,0) (2,0) (2,4) (0,4)} \closedcycle;
          \addplot[fill=uniAzul!40, draw=none, domain=0:2] {x^2} \closedcycle;
          \addplot[curva, domain=0:2.1] {x^2};
          \draw[dashed] (axis cs:2,0) -- (axis cs:2,4) -- (axis cs:0,4);
          \node at (axis cs:1.5,0.8) {$\tfrac13$};
          \node at (axis cs:0.7,2.6) {$\tfrac23$};
        \end{axis}
      \end{tikzpicture}
    \end{column}
    \begin{column}{0.5\textwidth}
      \begin{itemize}
        \item Área bajo la parábola: $\tfrac13 b\cdot b^2$.
        \item Área entre la parábola y el eje $y$: $\tfrac23 b^3$.
      \end{itemize}
      \begin{ejemplo}[análogo]
        Con $\displaystyle\sum_{i=1}^n i=\frac{n(n+1)}{2}$ se prueba que
        $\displaystyle\int_0^b x\dx=\frac{b^2}{2}$ (área de un triángulo), y con
        $\displaystyle\sum_{i=1}^n i^3=\left[\frac{n(n+1)}{2}\right]^2$ que
        $\displaystyle\int_0^b x^3\dx=\frac{b^4}{4}$.
      \end{ejemplo}
    \end{column}
  \end{columns}
\end{frame}

\begin{frame}{Sumas de Riemann (adelanto)}
  \begin{definicion}
    Si en cada $I_i$ elegimos un punto $c_i$ (punto de muestra), la \resaltar{suma de Riemann} es
    \[ S(f,P,\{c_i\})=\sum_{i=1}^n f(c_i)\,\Delta x_i. \]
  \end{definicion}
  Como $m_i\le f(c_i)\le M_i$, siempre $L(f,P)\le S(f,P,\{c_i\})\le U(f,P)$. Si $f$ es integrable:
  \[ \int_a^bf(x)\dx=\lim_{\|P\|\to0}\sum_{i=1}^nf(c_i)\Delta x_i. \]
  \begin{ejemplo}
    Con punto medio, $f(x)=x^2$ en $[0,2]$, $n=4$: \
    $S=\tfrac12\left[(0.25)^2+(0.75)^2+(1.25)^2+(1.75)^2\right]=2.625$.\quad (Exacto: $2.6\overline{6}$.)
  \end{ejemplo}
\end{frame}

\begin{frame}{Actividad colaborativa: Sumas de Riemann en una hoja de cálculo}
  \small
  \begin{activa}[Aprendizaje basado en la experiencia]
    Para $f(x)=x^2$ en $[0,1]$ calcule las sumas inferior y superior con $n=4,8,16,32$. \textbf{Conjeture} el valor común y la fórmula de las sumas en $n$; demuéstrela con $\sum k^2=\frac{n(n+1)(2n+1)}6$. Discuta por qué la función de Dirichlet no es integrable.
  \end{activa}
  \begin{itemize}
    \item \textbf{Organización:} equipos de 3 a 4 con roles rotativos (coordinador, redactor, verificador, portavoz); 12--15 minutos.
    \item \textbf{Producto:} Tabla de sumas, fórmula demostrada y argumento para Dirichlet.
    \item \textbf{Retroalimentación:} el docente recorre los equipos, hace preguntas y no da la solución hasta la puesta en común.
  \end{itemize}
\end{frame}

\section{Ejercicios}

\begin{frame}{Ejercicios propuestos}
  \begin{ejercicio}
    \begin{enumerate}
      \item Sea $f(x)=2x+1$ en $[0,2]$ y $P=\{0,\tfrac12,1,\tfrac32,2\}$. Calcule $L(f,P)$ y $U(f,P)$.
        Agregue el punto $\tfrac14$ y verifique el lema de refinamiento.
      \item Usando particiones regulares, calcule $\displaystyle\int_0^b x\dx$ y $\displaystyle\int_1^3(x^2+1)\dx$.
      \item Pruebe por definición que $\displaystyle\int_0^b x^3\dx=\frac{b^4}{4}$.
      \item Pruebe que $\displaystyle\int_a^b x^2\dx=\frac{b^3-a^3}{3}$ ($0\le a<b$) usando $x_i=a+i\frac{b-a}{n}$.
      \item Pruebe que si $f$ es constante, $f(x)=k$, entonces $\int_a^b f=k(b-a)$.
      \item Sea $f(x)=0$ si $x\neq\tfrac12$ y $f(\tfrac12)=1$ en $[0,1]$. Pruebe que $f$ es integrable y $\int_0^1f=0$.
      \item ¿Es integrable en $[0,1]$ la función $g(x)=x$ si $x\in\Q$ y $g(x)=0$ si $x\notin\Q$? Calcule
        $\underline{\int_0^1}g$ y $\overline{\int_0^1}g$.
      \item Pruebe que $f(x)=e^x$ es integrable en $[0,1]$ y calcule $\int_0^1e^x\dx$ usando
        $\sum_{i=0}^{n-1}r^i=\frac{r^n-1}{r-1}$.
    \end{enumerate}
  \end{ejercicio}
\end{frame}

\begin{frame}{Metacognición y uso responsable de la IA}
  \small
  \begin{columns}[t]
    \begin{column}{0.5\textwidth}
      \begin{metacog}[cierre, 3 min]
        \begin{enumerate}
          \metacomun
          \item ¿Puedo explicar en palabras por qué el refinamiento acerca las sumas inferior y superior?
        \end{enumerate}
      \end{metacog}
    \end{column}
    \begin{column}{0.5\textwidth}
      \begin{iabox}
        Pida a una IA la suma superior de $x^2$ con $n$ subintervalos. Verifique la fórmula con $n=1$ y $n=2$.\\[2pt]
        \textit{Regla:} primero intente usted; luego use la herramienta para contrastar, nunca para sustituir su razonamiento. Anote el cambio en su bitácora.
      \end{iabox}
    \end{column}
  \end{columns}
\end{frame}

\begin{frame}{Resumen}
  \begin{itemize}
    \item Partición $P$, norma $\|P\|$; refinamiento $P\subset P^*$.
    \item Sumas de Darboux: $L(f,P)=\sum m_i\Delta x_i$, \ $U(f,P)=\sum M_i\Delta x_i$.
    \item Refinar aumenta $L$ y disminuye $U$; toda $L$ es menor que toda $U$.
    \item $\displaystyle\underline{\int_a^b}f=\sup_PL$, \ $\displaystyle\overline{\int_a^b}f=\inf_PU$;
      integrable si coinciden. Criterio: $U-L<\varepsilon$.
    \item $\displaystyle\int_0^bx^2\dx=\frac{b^3}{3}$.
  \end{itemize}
  \vspace{3mm}
  \textbf{Próxima sesión:} propiedades de la integral definida y funciones integrables.
\end{frame}


\begin{frame}{Evaluación y trabajo posterior}
  \small
  \begin{evalbox}[Práctica del miércoles de la semana 2: Dirigida 1]
    \textbf{Cobertura:} semana 1. \textbf{Naturaleza:} práctica dirigida en equipo, con retroalimentación inmediata.
  \end{evalbox}
  \begin{itemize}
    \item \criterios
    \item \declaracionia
    \item \textbf{Portafolio:} guarde el producto de la actividad, la bitácora de metacognición y los ejercicios propuestos de esta sesión; son evidencia de progreso y base de la retroalimentación.
  \end{itemize}
\end{frame}

\end{document}
